\subsection*{S16. Heterogeneous array}

A single array whose elements have types with no common supertype.

\begin{lstlisting}[style=json]
{ "id": 1, "mixed": [ 1, "two", { "three": 3 } ] }
\end{lstlisting}

This is the first scenario where the widening used everywhere above stops
working. S1 and S6 widen along a chain --- $\bot \sqsubseteq \texttt{int}
\sqsubseteq \texttt{float}$ --- so any two observations have a least upper
bound. An integer, a string and an object have none, so there is nothing to
widen to.

\options
\begin{enumerate}[label=\textbf{(\alph*)}]
  \item \textbf{Reject}, with an error naming the path and the conflicting
        types. \selected
  \item \textbf{A generated variant.} The child table becomes a sparse table
        --- \texttt{mixed(id, parent\_id, idx, tag, v\_int, v\_str,
        v\_obj\_id)} with exactly one value column non-null per row --- and
        the element type becomes
        \texttt{type mixed\_elt = Int of int | Str of string | Obj of
        Mixed\_obj.id}.
  \item \textbf{An opaque raw-JSON column}, leaving elements as
        \texttt{Yojson.Safe.t} for the consumer to interpret.
  \item \textbf{Widen to string}, serializing non-string elements back to JSON
        text.
\end{enumerate}

\decision{Reject. An array whose elements have no least upper bound is an
error naming the path and the types observed.}

\paragraph{What does not count as heterogeneous.}
The rejection is narrow, and several mixtures that look heterogeneous are
resolved by decisions already taken:

\begin{itemize}
  \item \texttt{[1, 2.5]} --- integers and floats widen to \texttt{float}
        under S1a. Not a conflict.
  \item \texttt{[\{"a":1\}, \{"b":2\}]} --- objects with different members
        merge into one table whose columns are partial under S4. Not a
        conflict.
  \item \texttt{[[1], [2]]} --- arrays of arrays are S17, not this scenario.
\end{itemize}

\noindent
What is rejected is a genuine absence of a least upper bound: integer against
string, integer against boolean, string against boolean, or any scalar against
an object or an array.

\rationale{Option (b) is the principled answer and the one a relational
modeller would reach for, but it reintroduces the data-dependence problem of
S9 in a third costume: the \emph{shape of the generated type} would depend on
which element types happened to appear, so adding a document could add a
constructor and change the generated variant. Option (d) discards the type
information it claims to preserve, and option (c) defers the whole problem to
the consumer while making the column untyped.

Rejection also matches what such arrays usually are in practice. A corpus
mixing scalars and objects at one path generally reflects a fault in whatever
produced it, and is something the user would rather be told about.}

\limitation{Option (b) remains the natural extension if heterogeneous arrays
turn out to matter, and adopting it later would not disturb any other decision
in this document. Note that S18 asks the same question across documents rather
than within an array, and the two answers should be kept consistent.}

\subsection*{S16b. Null elements in an array}

S4 governs object members rather than array elements, so it does not by itself
decide the status of a null sitting in element position.

\begin{lstlisting}[style=json]
{ "mixed": [ 1, null, 2 ] }
\end{lstlisting}

\options
\begin{enumerate}[label=\textbf{(\alph*)}]
  \item \textbf{\texttt{null} is the absence of a value.} The element column
        becomes nullable, the element type is \texttt{int option}, and the
        array is not heterogeneous. \selected
  \item \textbf{\texttt{null} is a distinct type} with no least upper bound
        against \texttt{int}, making the array heterogeneous and rejecting it
        under S16.
\end{enumerate}

\decision{Treat \texttt{null} in element position as the absence of a value.
The element column is nullable and the element type gains an \texttt{option};
the array is not heterogeneous.}

\begin{lstlisting}[style=json]
{ "mixed": [ 1, null, 2 ] }
  -> element type int option; the value column is nullable

{ "mixed": [ null, null ] }
  -> element type never observed; S15 applies, with a nullable column
\end{lstlisting}

\rationale{Consistent with the rest of this document, where \texttt{null} has
uniformly meant the absence of a value rather than a type of its own. S6
already gave an all-null path a home rather than treating it as a conflict,
and rejecting \texttt{[1, null, 2]} while accepting an all-null scalar field
would be hard to justify.

The contrary reading was weighed: an array has no notion of a missing element,
so \texttt{[1, null, 2]} has three elements and the middle one is a value that
happens to be null, which makes \texttt{null} more type-like in element
position than in member position. It is recorded here and not adopted.}

\paragraph{A collapse that does not occur.}
Unlike S4, this introduces no loss. In member position \texttt{None} is
ambiguous between an explicit \texttt{null} and an absent key. In element
position there is no absent case to collide with --- every index up to the
array's length is occupied --- so a null element reconstructs as \texttt{null}
unambiguously. Element nullability is lossless where member nullability is
not.
